Gedhéné himpunan winates~$A$ bisa diukur nganggo wilangan asli.
Kita bisa mbandhingaké gedhéné himpunan winates kanthi mbandhingaké
wilangan-wilangan mau. Kanggo himpunan tanpa wates, carané luwih
ruwet. Tataran kapisan yaiku himpunan \emph{tanpa wates kang bisa
dietung}. Himpunan~$A$ \emph{bisa dietung} yèn anggotané bisa
ditata ing \emph{enumerasi}, yaiku dhaptar tanpa wates ing siji
arah. Iki kelakon yèn ana fungsi surjèktif $f\colon\PosInt\to A$.
Himpunan iku tanpa wates kang bisa dietung yèn bisa dietung lan
ora winates. \emph{Cara zèg-zag Cantor} nuduhaké yèn himpunan
pasangan anggota saka himpunan tanpa wates kang bisa dietung
uga bisa dietung. Cara iku uga nuduhaké yèn $\Rat$, himpunan
wilangan rasional, bisa dietung.

Ana uga himpunan tanpa wates kang \emph{ora bisa dietung}.
Iki bisa dibuktèkaké nganggo \emph{argumèn diagonal} utawa
\emph{reduksi}. Ing argumèn diagonal, upamakna himpunan~$A$
bisa dietung. Saka enumerasi hipotètis iku, tetepna sawijining
anggota~$A$ kanthi cara kang njamin anggota iku béda karo saben
anggota ing dhaptar. Mula dhaptar iku ora ngemot kabèh anggota~$A$.
Ora ana enumerasi kang bisa ngemot kabèh anggota~$A$.
Ing reduksi, kanggo nuduhaké yèn $A$ ora bisa dietung, pasangna
saben anggota~$A$ marang anggota sawijining himpunan~$B$ kang wis
dikawruhi ora bisa dietung, kanthi surjèktif. Enumerasi hipotètis
saka~$A$ banjur bakal ngasilaké enumerasi saka~$B$. Amarga $B$
ora bisa dietung, enumerasi~$A$ ora ana.

Sacara umum, gedhéné himpunan tanpa wates bisa dibandhingaké.
$A$ lan $B$ \emph{padha cacahé} yèn ana bijèksi ing antarané.
$A$ ora luwih gedhé tinimbang~$B$, $\cardle{A}{B}$, yèn ana
fungsi injèktif saka $A$ menyang~$B$. Miturut teorema
Schr\"oder--Bernstein, iki mènèhi urutan miturut gedhéné
himpunan tanpa wates. \emph{Teorema Cantor} nyatakaké yèn
kanggo saben $A$, $\cardless{A}{\Pow{A}}$. Iki ngambakaké
asil yèn $\Pow{\PosInt}$ ora bisa dietung lan nuduhaké yèn
tataran tanpa wates cacahé uga tanpa wates.
